\scanpage{19}

dans la suite des $\operatorname{div}(n_{i\alpha})$ est $f_{i\ell}$, et $\Phi(i,\ell)_k$ est le \deleted{complexe} \inserted{module $\frac12$ \unsure{simpl.|simplicial}} associé à

\diagram{\begin{tikzcd}[ampersand replacement=\&,column sep=small]
{}-{} \& 0\arrow[r] \& \overset{\mathrm{deg}\,i}{k}\arrow[r,"n"] \& \overset{\mathrm{deg}\,i-1}{k}\arrow[r] \& 0-
\end{tikzcd}}

où $\operatorname{div}(n)=\ell$. On a donc une extension

\diagram{\begin{tikzcd}[ampersand replacement=\&,column sep=small]
0\arrow[r] \& \Lambda^{i-1}\Psi\arrow[r] \& \Phi(i,n)\arrow[r,dash] \& \Lambda^i\Psi\arrow[r] \& 0
\end{tikzcd}}

où l'hom associé à cette extension

\[ k\longrightarrow k \]

(cf ---) est $n\,\mathrm{id}$, où $\operatorname{div}(n)=\ell$, donc

\[
\deleted{\unsure{H_j}}\pi_j(\Phi(i,n))=
\begin{cases}0&\text{si }j\ne i-1\\
k/\operatorname{div}n&\text{si }j=i-1.
\end{cases}
\]

\uline{Remarques} 1) La fonction de Hilbert de $\Phi(i,n)$ est donc la $=$ que celle de $\Lambda^i\Phi_{\bullet}$, savoir $\binom{n+1}{i}$. Donc si $L_{\bullet}$ est donné par le produit de la prop. ---.

\[
\operatorname{rg}L_n=\sum_{i\leq n+1}(\mu_i+\beta_i)\binom{n+1}{i}
+\sum_{0\leq j\leq n}\nu_j\binom{n}{i}.
\]

\uline{Donc} \unsure{1$'$)|h$'$)|2)} $L_{\bullet}$ \deleted{\unsure{en}} est \inserted{$\partial$-parfait, il est} connu à isom (non unique) près quand on connaît a) $\operatorname{rg}\pi_i(A_{\bullet})=\nu_i$ et le nombre $\beta_{i+1}$ de ses diviseurs de torsion, pour $\forall\,i\geq0$, et b) la fonction $f_{L_{\bullet}}$ \unsure{res.|res|très} pour $n$ grand (qui coïncide \deleted{alors} avec \inserted{pour $n$ grand} le polynôme $\sum_{i\geq1}(\mu_i+f_i)\binom{n+1}{i}+\sum_{j\geq0}\nu_j\binom{n}{i}$

% archive sheet number: 18
